% ======================================================================
% common.tex -- shared preamble for "Probability From the Ground Up"
% Visual language matched to the companion ground-up readers.
% ======================================================================
\usepackage[margin=0.95in]{geometry}
\usepackage{amsmath,amssymb,amsthm}
\usepackage{mathtools}
\usepackage{enumitem}
\usepackage{longtable}
\usepackage{array}
\usepackage{xcolor}
\usepackage{hyperref}
\usepackage{fancyhdr}
\usepackage{titlesec}
\usepackage{tcolorbox}
\tcbuselibrary{skins,breakable}
\usepackage{multicol}
\usepackage{booktabs}
\usepackage{needspace}
\usepackage{float}
\usepackage{listings}
\usepackage{tikz}
\usetikzlibrary{arrows.meta,positioning,calc,fit,shapes.geometric,shapes.misc,decorations.pathreplacing,backgrounds,matrix,chains,trees,patterns,angles,quotes,intersections}
\usepackage{pgfplots}
\pgfplotsset{compat=1.16}
\usepgfplotslibrary{fillbetween,statistics}

\hypersetup{colorlinks=true, linkcolor=blue!50!black, urlcolor=blue!50!black, citecolor=blue!50!black}
\setlength{\parskip}{0.62em}
\setlength{\parindent}{0pt}
\setlength{\headheight}{15pt}
\setlist[itemize]{topsep=0.2em,itemsep=0.2em}
\setlist[enumerate]{topsep=0.2em,itemsep=0.25em}

\pagestyle{fancy}
\fancyhf{}
\rhead{\thepage}
\lhead{\leftmark}
\renewcommand{\headrulewidth}{0.4pt}

% ---- Conversational boxes (shared visual language) -------------------
\newtcolorbox{focusbox}{colback=blue!3!white,colframe=blue!50!black,title=Focus,breakable}
\newtcolorbox{intuitionbox}{colback=green!3!white,colframe=green!40!black,title=Intuition,breakable}
\newtcolorbox{formalbox}{colback=gray!4!white,colframe=gray!70!black,title=Formal,breakable}
\newtcolorbox{walkbox}{colback=orange!5!white,colframe=orange!60!black,title=Walking question,breakable}
\newtcolorbox{trapbox}{colback=red!3!white,colframe=red!55!black,title=Common trap,breakable}
\newtcolorbox{historybox}{colback=yellow!7!white,colframe=yellow!45!black,title=History,breakable}
\newtcolorbox{happeningbox}{colback=green!5!white,colframe=green!55!black,title=What is really happening,breakable}
\newtcolorbox{paradoxbox}{colback=red!4!white,colframe=red!60!black,title=Paradox,breakable}
\newtcolorbox{labbox}{colback=purple!3!white,colframe=purple!45!black,title=Work it yourself,breakable}
\newtcolorbox{aestheticbox}{colback=cyan!4!white,colframe=cyan!50!black,title=Why this matters,breakable}
\newtcolorbox{applicationbox}{colback=magenta!4!white,colframe=magenta!55!black,title=Application,breakable}
\newtcolorbox{connectionbox}{colback=teal!4!white,colframe=teal!55!black,title=Connection,breakable}
\newtcolorbox{whybox}{colback=blue!4!white,colframe=blue!60!black,title=Why it works,breakable}
\newtcolorbox{marginnotebox}{colback=gray!6!white,colframe=gray!55!black,title=Reading note,breakable}
\newtcolorbox{recoverybox}{colback=white,colframe=gray!45!black,title=Pause,breakable}

% ---- Colors used by diagrams ----------------------------------------
\definecolor{sourceblue}{RGB}{61,105,180}
\definecolor{targetgreen}{RGB}{62,145,95}
\definecolor{quotientorange}{RGB}{214,127,35}
\definecolor{trapred}{RGB}{190,65,55}
\definecolor{structurepurple}{RGB}{120,88,165}
\definecolor{distblue}{RGB}{40,90,170}
\definecolor{histfill}{RGB}{120,170,220}
\definecolor{densityred}{RGB}{200,60,50}
\definecolor{codebg}{RGB}{247,247,250}

% ---- Code listings ---------------------------------------------------
\lstdefinestyle{book}{
  basicstyle=\ttfamily\small, backgroundcolor=\color{codebg},
  keywordstyle=\color{sourceblue}\bfseries, commentstyle=\color{targetgreen}\itshape,
  stringstyle=\color{quotientorange}, numberstyle=\tiny\color{gray},
  numbers=left, numbersep=8pt, frame=single, framesep=5pt, rulecolor=\color{gray!40},
  breaklines=true, showstringspaces=false, columns=fullflexible, keepspaces=true, tabsize=2,
}
\lstset{style=book}

% ---- Math shorthands -------------------------------------------------
\newcommand{\N}{\mathbb{N}}
\newcommand{\Z}{\mathbb{Z}}
\newcommand{\Q}{\mathbb{Q}}
\newcommand{\R}{\mathbb{R}}
\newcommand{\C}{\mathbb{C}}
\newcommand{\eps}{\varepsilon}
\newcommand{\dd}{\mathrm{d}}
\renewcommand{\P}{\operatorname{\mathbb{P}}}
\newcommand{\E}{\operatorname{\mathbb{E}}}
\newcommand{\Var}{\operatorname{Var}}
\newcommand{\Cov}{\operatorname{Cov}}
\newcommand{\Corr}{\operatorname{Corr}}
\newcommand{\SD}{\operatorname{SD}}
\newcommand{\ind}{\mathbf{1}}
\newcommand{\given}{\mid}
\newcommand{\indep}{\perp\!\!\!\perp}
\newcommand{\dto}{\xrightarrow{d}}
\newcommand{\pto}{\xrightarrow{p}}
\newcommand{\asto}{\xrightarrow{a.s.}}
\newcommand{\Normal}{\mathcal{N}}
\newcommand{\Bern}{\operatorname{Bernoulli}}
\newcommand{\Bin}{\operatorname{Binomial}}
\newcommand{\Pois}{\operatorname{Poisson}}
\newcommand{\Geom}{\operatorname{Geometric}}
\newcommand{\Unif}{\operatorname{Uniform}}
\newcommand{\Expo}{\operatorname{Exponential}}
\newcommand{\Gam}{\operatorname{Gamma}}
\newcommand{\Bet}{\operatorname{Beta}}
\newcommand{\KL}{\operatorname{KL}}
\DeclarePairedDelimiter{\abs}{\lvert}{\rvert}
\DeclarePairedDelimiter{\norm}{\lVert}{\rVert}
\DeclarePairedDelimiter{\set}{\{}{\}}
\DeclarePairedDelimiter{\floor}{\lfloor}{\rfloor}
\DeclarePairedDelimiter{\ceil}{\lceil}{\rceil}

% ---- Theorem-like environments --------------------------------------
\theoremstyle{definition}
\newtheorem{definition}{Definition}[chapter]
\newtheorem{example}{Example}[chapter]
\newtheorem{exercise}{Exercise}[chapter]
\theoremstyle{plain}
\newtheorem{theorem}{Theorem}[chapter]
\newtheorem{lemma}{Lemma}[chapter]
\newtheorem{corollary}{Corollary}[chapter]
\newtheorem{proposition}{Proposition}[chapter]
\theoremstyle{remark}
\newtheorem{remark}{Remark}[chapter]

\newcommand{\chapterline}{\vspace{0.4em}\hrule\vspace{0.8em}}
